\section{OBJETIVOS}
\justifying

Con el fin de lograr un sistema óptimo que cumpla con un correcto funcionamiento y fiabilidad, se plantearon objetivos para establecer el rumbo del proyecto, definiendo de forma clara y precisa lo que se pretende alcanzar.

\subsection{Objetivo General}
\begin{itemize}[noitemsep]
    \item Diseñar un sistema didáctico modular para la experimentación de electrónica de potencia, con el fin de facilitar el aprendizaje práctico mediante el uso de módulos especializados y material teórico complementario; enfatizando la simplicidad de uso, seguridad y adaptabilidad para futuros desarrollos.
\end{itemize}

\subsection{Objetivos Específicos}
% Pueder ser mejor usar items y no enumeración
\begin{enumerate}[noitemsep]
    \item Seleccionar los módulos a implementar en el sistema didáctico.
    \item Desarrollar los circuitos y placas de los módulos, considerando aspectos de potencia y disipación de calor.
    \item Diseñar los módulos con la intensión de ser fácilmente conectables e intercambiables.
    \item Seleccionar los componentes y materiales necesarios para la construcción de los módulos.
    \item Diseñar los gabinetes y la estructura del sistema para facilitar su impresión en filamento 3D.
    \item Adaptar el sistema para futuras expansiones y mejoras.
    \item Validar el correcto funcionamiento de los módulos mediante pruebas y ajustes.
    \item Diseñar una interfaz de usuario intuitiva y accesible.
    \item Asegurar la seguridad del sistema y de los usuarios.
    \item Evaluar el funcionamiento del sistema en su totalidad.
    \item Elaborar un manual de usuario que sirva como complemento teórico y facilite la puesta en marcha del sistema.
    \item Garantizar el fácil mantenimiento y reparación del sistema, en caso de ser necesario.    
    \item Validar el funcionamiento del sistema a través de pruebas con usuarios finales.
\end{enumerate}